%% ============================================================
%% 分片 27：原文第 10 节「The Bogomolny Equations And The Space Of Hecke Modifications」（原文 glshort.tex 起于第 8120 行）
%% 由 tools/harvmac2tex.py 从 arXiv 源码（hep-th/0604151v3）自动生成。
%% 只译英文散文；公式、编号、\label、\cite 一律不改。
%% ============================================================

\subsection{单极子冒泡} \label{sec:monbub}

在第~\ref{sec:spaceheck}~节中我们发现，对某个给定类型
$^L\negthinspace w$ 的 Hecke 修正空间
${{\EUY}}(^L\negthinspace w)$，其自然紧化
$\bar {{\EUY}}(^L\negthinspace w)$ 是这样实现的：把若干相伴类型的、
权重更小的 Hecke 修正一并包含进来。

现在我们就能从 Bogomolny 方程的角度看出这一点的缘由。考虑
$\Bbb{R}^3$ 上在原点处具有指定奇异性的一个 Bogomolny 方程解。它
对应于 $\Bbb{C}^2\cong \Bbb{R}^4$ 上的一个等变瞬子。$\Bbb{R}^4$
上的瞬子方程是尺度不变的，因此会出现数学文献中称为「冒泡」
（bubbling）的现象：一个瞬子可以收缩成一点。对于
${\EUF}$-等变瞬子，这个点必然就是 $\Bbb{C}^2$ 的原点。

借助等变瞬子与三维中具有点奇异性的 Bogomolny 方程解之间的
关系 \cite{ref110}（我们把这种解解释为 't Hooft 算子），把上述
图像放到三维来看，就得到一种或可称为「单极子冒泡」（monopole
bubbling）的现象。单极子冒泡（Kronheimer 在未发表的工作中也
独立发现了它，与 \cite{ref110} 有关）之所以相对陌生，是因为它
对于 $\Bbb{R}^3$（或其他完备黎曼流形）上满足通常边界条件的光滑
Bogomolny 方程解并不会发生。恰恰相反，$\Bbb{R}^3$ 上的光滑
单极子具有离散的拓扑分类与特征尺寸，后者涉及 $\Bbb{R}^3$ 上
无穷远处的渐近行为。单极子冒泡只能出现在有 't Hooft 算子存在的
情形；它是一个过程：一个量子化的电荷单位集中到一点（即
't Hooft 算子所在的位置）并随之消失。

下面我们非常简略地利用 $\Bbb{R}^4\cong \Bbb{C}^2$ 上瞬子的
ADHM 构造（ADHM construction）\cite{ref125}，来确定当规范群为
$G=SU(N)$ 时，等变瞬子的冒泡——换言之，单极子冒泡——如何影响
't Hooft 算子的权 $^L\negthinspace w$。（对于其他存在 ADHM
构造的群，即经典群 $SO(N)$ 与 $Sp(N)$，也可以作类似的分析。但
例外群没有 ADHM 构造。）我们将看到，冒泡对权的改变恰好与预期
相符。

我们将沿用 \cite{ref126} 第 3.3 节。为描述 $\Bbb{R}^4$ 上的
$SU(N)$ 瞬子，设 $U$ 是一个二维向量空间（它作为 $\Bbb{R}^4$ 的
旋量丛之一出现），${\EUF}$ 在它上面按矩阵 % 式 (10.6)
\setcounter{equation}{5}
\begin{equation}\label{eq:normalo}
\left(\begin{matrix}e^{i\theta}
& 0 \\ 0 & e^{-i\theta}\\\end{matrix}\right).
\end{equation}
作用。对任何带有 ${\EUF}$ 作用的向量空间 $V$，记 $\chi(V)$ 为
${\EUF}$ 的特征；例如
$\chi(U)=e^{i\theta}+e^{-i\theta}$。为构造瞬子数为 $k$ 的瞬子，
设 ${{\EUH}}$ 是一个 $k$ 维向量空间。再设 $E_\infty$ 是一个
秩为 $N$ 的向量空间（它实际上是瞬子丛在无穷远处的纤维）。瞬子
由如下一个特殊类型的复形
% 式 (10.7)
\setcounter{equation}{6}
\begin{equation}\label{eq:onnos}
0\,\underarrow{~~~~}\,{{\EUH}}\underarrow{\,\,\alpha(z)\,\,}{{\EUH}}\otimes
U\oplus E_\infty\underarrow{\,\,\beta(z)\,\,}
{{\EUH}}\,\underarrow{~~~~}\, 0,
\end{equation}
 其中 $\beta(z)\alpha(z)=0$。就我们此处的有限目的而言，并不需要
这些细节。$\alpha(z)$ 与 $\beta(z)$ 线性依赖于 $\Bbb{C}^2$ 的
复坐标 $z_1,z_2$。对于光滑瞬子，该复形对任何
$z=(z_1,z_2)\in \Bbb{C}^2$ 唯一非零的上同调群是
$H^1(z)={\rm ker}\,\beta(z)/{\rm im}\,\alpha(z)$。由 ADHM 数据
得到的瞬子丛 $E$ 在 $z$ 处的纤维就是 $E_z=H^1(z)$。

为了得到 ${\EUF}$-等变瞬子，我们选取 ${\EUF}$ 在 $E_\infty$ 与
${{\EUH}}$ 上的作用，并要求复形 \eqref{eq:onnos}\ 是
${\EUF}$-等变的。瞬子丛在原点的纤维（记作 $E_0$）特别地就是
相应复形的第一上同调群：
% 式 (10.8)
\setcounter{equation}{7}
\begin{equation}\label{eq:otonos}
0\longrightarrow{{\EUH}}\underarrow{\,\,\alpha(0)\,\,}{{\EUH}}\otimes
U\oplus E_\infty\underarrow{\,\,\beta(0)\,\,}
{{\EUH}}\longrightarrow 0,
\end{equation}
 对一个仅在 $H^1$ 处有非零上同调的等变复形
$0\underarrow{~~~} V_0\underarrow{\,\,\alpha\,\,}
V_1\underarrow{\,\,\beta\,\,} V_2\underarrow{~~~} 0$，其 $H^1$ 的
特征可以写成
$\chi(H^1)=-\chi(V_0)+\chi(V_1)-\chi(V_2)$。特别地
% 式 (10.9)
\setcounter{equation}{8}
\begin{equation}\label{eq:onsno}
\chi(E_0)=\chi(E_\infty)+\chi({{\EUH}})(e^{i\theta}+e^{-i\theta}-2).
\end{equation}

$\chi(E_0)$ 之所以重要，是因为按照 Kronheimer 的构造，当把该
等变瞬子放到三维中解释时，它决定了出现在
$\Bbb{R}^3=\Bbb{C}^2/{\EUF}$ 原点的 't Hooft 算子的权。若
$\chi(E_0)=\sum_{i=1}^N\exp({im_i\theta})$，则这些权为
$^L\negthinspace w=(m_1,m_2,\dots,m_N)$。

当 $\alpha$ 与 $\beta $ 变动到如下情形时，瞬子冒泡就会发生：对
某个分解 ${{\EUH}}={{\EUH}}_1\oplus
{{\EUH}}_2$，复形 \eqref{eq:onnos}\ 分解为两个复形的直和——一个
零复形 % 式 (10.10)
\setcounter{equation}{9}
\begin{equation}\label{eq:teto}
0\,\underarrow{~~~~}\,
{{\EUH}}_1\underarrow{\,\,\,0\,\,\,}{{\EUH}}_1\otimes
U\underarrow{\,\,\,0\,\,\,}{{\EUH}}_1\,\underarrow{~~~~}\, 0
\end{equation}
 与一个复形
% 式 (10.11)
\setcounter{equation}{10}
\begin{equation}\label{eq:ponos}
0\,\underarrow{~~~~}\,{{\EUH}}_2\underarrow{\,\,\alpha(z)\,\,}{{\EUH}}_2\otimes
U\oplus E_\infty\underarrow{\,\,\beta(z)\,\,}
{{\EUH}}_2\,\underarrow{~~~~}\, 0
\end{equation}
 的直和，而后者的上同调只出现在 1 次。所谓「冒泡后的」瞬子丛
$E'$，其瞬子数比原来少了 ${{\EUH}}_1$ 的秩，它以这第二个复形的
第一上同调群为纤维；与 \eqref{eq:onsno}\ 类似，它在原点的纤维
的特征为
$\chi(E'_0)=\chi(E_\infty)+\chi({{\EUH}}_2)(e^{i\theta}+e^{-i\theta}-2)$。
于是
% 式 (10.12)
\setcounter{equation}{11}
\begin{equation}\label{eq:ponzo}
\chi(E'_0)=\chi(E_0)+\chi({{\EUH}}_1)(2-e^{i\theta}-e^{-i\theta}).
\end{equation}

这个公式决定了瞬子冒泡下权的改变。$E_0$ 的权 $m_i$ 与 $E'_0$ 的
权 $m'_i$ 之间满足
% 式 (10.13)
\setcounter{equation}{12}
\begin{equation}\label{eq:yhog}
\sum_{i=1}^N\exp({im'_i\theta})=
\sum_{i=1}^N\exp({im_i\theta})+\chi({{\EUH}}_1)(2-e^{i\theta}-e^{-i\theta}).
\end{equation}

例如，若 ${{\EUH}}_1$ 是一维的、特征为 $\exp(im\theta)$（$m$ 为
某个整数），则 $E_0$ 的权表必须包含 $m+1$ 与 $m-1$（因为
$\chi(E'_0)$ 是系数非负的 Fourier 级数）。\eqref{eq:yhog}\ 表明，
在这次冒泡过程中恰好有两个权发生变化，即
$(m+1,m-1)\to (m,m)$。这对应于我们在第~\ref{sec:spaceheck}~节
讨论过的 't Hooft/Hecke 算子的最简单退化。再举一例：设
${{\EUH}}_1$ 的秩为二，特征为 $\exp({im
\theta})+\exp({im'\theta})$。若 $|m-m'|>1$，我们只是得到上述
类型的一对移动，即把权 $(m+1,m-1,m'+1,m'-1)$ 换成
$(m,m,m',m')$。但若 $m'=m+1$，则 \eqref{eq:yhog}\ 还允许另一种
移动：把权 $(m+2,m-1)$ 换成 $(m+1,m)$。这对应于另一种预期的
't Hooft/Hecke 算子退化。更一般地，若 ${{\EUH}}_1$ 的秩为 $s$、
权为 $m,m+1,\dots,m+s-1$，则 \eqref{eq:yhog}\ 导致一次移动%
\footnote{要实现这次移动，取 ${{\EUH}}={{\EUH}}_1$、
${{\EUH}}_2=0$，并令 $E'$ 为以 $E_\infty$ 为纤维的平凡丛，
即可得到一个具有适当 ADHM 形式的复形。在正文已给出 $E'_0$ 与
${{\EUH}}_1$ 的特征的情况下，满足 \cite{ref126}\ 的式 (3.3.11)
并且给出一个光滑的、瞬子数为 $s$、能通过冒泡退化为 $E'$ 的
瞬子丛 $E$ 的 $\tau_1,\tau_2,\sigma,$ 与 $\pi$ 的
${\EUF}$-不变选取，本质上是唯一的。}：其中 $E_0$ 的权
$(m+s,m-1)$ 被换成 $E'_0$ 的权 $(m+s-1,m)$。对于 $G=SU(N)$，
't Hooft/Hecke 算子的预期退化全都是这种类型移动的复合。对于
一般的 $G$，预期的退化可以由嵌入 $G$ 的各个子群中的 $SU(2)$
瞬子的冒泡生成。

特别要注意：在不要求 ${\EUF}$-不变时，可以认为瞬子是一个一个地
冒泡的；但一旦施加 ${\EUF}$-不变性，情况就不再如此。存在一些
不可约的过程，其中任意多个 ${\EUF}$-不变瞬子同时冒泡。关于
$\Bbb{R}^3$ 上某些 Bogomolny 模空间的详细描述（在适当的极限下
可以从中看到单极子冒泡现象），参见 \cite{ref112,ref113}。

\par\medskip{\it 奇点处的超凯勒结构}

在第~\ref{sec:spaceheck}~节中，我们曾用一个例子指出：给定类型
$^L\negthinspace w$ 的 Hecke 修正的紧化空间
$\bar {{\EUY}}(^L\negthinspace w)$ 虽然不是超凯勒流形，但它的
奇点确实具有超凯勒结构。现在我们可以解释其中的缘由。这些奇点
来自单极子冒泡，而单极子冒泡由 $\Bbb{C}^2=\Bbb{R}^4$ 上
${\EUF}$-等变瞬子的模空间所支配。但 $\Bbb{R}^4$ 上的瞬子模空间
从 $\Bbb{R}^4$ 的超凯勒结构继承了一个超凯勒结构（这一点在 ADHM
构造中是显然的）。由于群 ${\EUF}$ 保持 $\Bbb{R}^4$ 的超凯勒
结构，${\EUF}$-等变瞬子的空间也是超凯勒的（同样，ADHM 构造使
这一点显而易见）。

我们也可以完全在三维的语言中论证——不动用与四维等变瞬子的
联系——单极子冒泡、从而 $\bar {{\EUZ}}({}^L\negthinspace
w_1,\dots,{}^L\negthinspace w_n)$ 的奇点，是由一个超凯勒模空间
支配的。诚然，$\Bbb{R}^3$ 上 Bogomolny 方程解的模空间是超凯勒的
\cite{ref109}。\footnote{这可以用我们稍后为说明 $W=I\times C$ 上
Bogomolny 方程解的模空间具有 Kähler 结构而使用的那类论证来证明。
那些论证在把 $I$ 换成 $\Bbb{R}$ 时同样有效。如果再取
$C=\Bbb{R}^2$，则 $W=\Bbb{R}^3$ 可以有许多种方式写成
$\Bbb{R}^2\times \Bbb{R}$，于是得到的是超凯勒结构，而不只是我们
在第~\ref{sec:kahstr}~节得到的那种 Kähler 结构。} 而
$\bar {{\EUZ}}({}^L\negthinspace w_1,\dots,{}^L\negthinspace
w_n)$ 的空间并不来自 $\Bbb{R}^3$ 上 Bogomolny 方程的解，而是来自
更复杂的三维流形 $W=I \times C$，并带有相当特别的边界条件。所以
$\bar {{\EUZ}}({}^L\negthinspace w_1,\dots,{}^L\negthinspace
w_n)$ 不是超凯勒的。然而，单极子冒泡是一个局部现象。当一个
单极子电荷单位收缩到一点时，其局部结构并不依赖于发生这一过程的
三维流形的细节。因此奇点附近的结构是超凯勒的，就如同身处
$\Bbb{R}^3$ 上一样。

关于这些超凯勒奇点的形变与消解，我们马上还要给出更多细节。

\begin{yizhu}[易错点]
按：本片反复出现的 $^L\negthinspace w$ 中，左上角的 $L$ 表示这是
Langlands 对偶群 ${}^LG$ 的权，而不是某个普通上标；负细空格命令只是
排版用的，不含数学内容。${\EUF}$ 是 macros.tex 里为绕原点的相位旋转
群 $U(1)$ 预留的花体字母宏（harvmac 习惯用 ${\EUS}$、${\EUF}$ 之类表示
「特殊」群），本片中「${\EUF}$-等变」即「$U(1)$-等变」，不要误当作
与 $F$ 有关的对象。另外，$G=SU(N)$ 的对偶群是 $PSU(N)$，所以权
$(m_1,\dots,m_N)$ 只定义到给所有 $m_i$ 同时加同一个常数（差一个中心
特征）；(10.13) 两边作为权表可以整体平移而不改变等式，正文比较
Fourier 系数时用的正是这一自由度。
\end{yizhu}

\begin{yizhu}[含糊点拨]
按：$\chi(H^1)=-\chi(V_0)+\chi(V_1)-\chi(V_2)$ 其实是三项复形的
Euler 特征 $\sum_i(-1)^i\chi(H^i)$；这里只有 $H^1$ 非零，故二者相等。
要点在于右端只依赖复形的项 $V_i$，与微分 $\alpha,\beta$ 无关；因此
即使 $H^1(z)$ 随 $z$ 变动而发生跳跃（这正是冒泡时发生的事），它的
特征仍可由 $E_\infty$ 与 ${{\EUH}}$ 的特征算出。原文「对光滑瞬子只有
$H^1$ 非零」一句省略了 ADHM 数据的非奇异性条件（即实瞬时方程与
稳定性条件，使 $\ker\alpha$ 与 $\mathrm{coker}\,\beta$ 不出现）；缺了它，复形
可以有 $H^0=\ker\alpha$ 或 $H^2$ 非零。
\end{yizhu}

\begin{yizhu}[物理图像]
按：「冒泡」的画面可以这样读：$\Bbb{C}^2=\Bbb{R}^4$ 上的等变瞬子有
一个被 $U(1)$ 对称性保住的尺度模，当它缩到原点时，三维图上看到的
就是「一个量子化的单极子电荷单位在 't Hooft 算子所在点被吞掉」。在
Hecke 修正一侧，这对应权表的一次退化 $(m+1,m-1)\to(m,m)$：两个权并成
一个双重权。换言之，Hecke 修正空间的紧化边界正是由这些冒泡构型铺
出来的；$\bar {{\EUY}}$ 的奇点就是气泡收缩的那一瞬间。
\end{yizhu}

\begin{yizhu}[易错点]
按：(10.12) 看似「凭空多出」一些权，但 $E'_0$ 是真实瞬子丛在原点的
纤维，其特征 $\sum_i\exp(im'_i\theta)$ 的系数必须是非负整数，因此
$\chi(E_0)+\chi({{\EUH}}_1)(2-e^{i\theta}-e^{-i\theta})$ 必须能写成
系数非负的 Fourier 级数；这正是正文用来反推「$E_0$ 必须含有 $m\pm1$」
的依据。注意 $2-e^{i\theta}-e^{-i\theta}$ 的系数 $(1,-1,-1)$ 并非全正，
所以冒泡只会朝「权更集中」的方向发生，不能从虚空中产生权。
\end{yizhu}

\begin{yizhu}[物理图像]
按：奇点为什么是超凯勒的？ADHM 构造把 $\Bbb{R}^4$ 上瞬子模空间写成
超凯勒商（三个 moment map 同时取零再取商），而 ${\EUF}$ 保持
$\Bbb{R}^4$ 的超凯勒结构，故等变瞬子空间继承超凯勒结构；单极子冒泡
由这个空间支配，因此 $\bar {{\EUY}}$ 的奇点局部是超凯勒的。
$\bar {{\EUY}}$ 整体反而不是超凯勒的，因为它来自带边界条件的三维
流形 $W=I\times C$ 上的 Bogomolny 解，而冒泡只是局部现象——原文最后
一段正是这个论证。
\end{yizhu}
